\subsection{Sparsity-independent one-level expander decomposition}
\label{sec:techoverviewub}
Given a graph $G=(V,E)$ in a  dynamic stream, and a parameter $\phi \in (0,1)$, we consider the problem of computing an \smash{$(\epsilon, \phi)$-ED} of $G$ for \smash{$\epsilon = O(\phi   \log n)$}. We show that one can do this in \smash{$\otil(n)$} bits of space, without any dependence on $\phi$.
In this section, we sketch our approach.

Let us begin by recalling the standard recursive framework for constructing expander decompositions~\cite{kvv,Trevisan08,stspectral}, concisely summarized in \cref{alg:metaed}. For any parameter $\phi \le 10^{-1}/\log n$, this algorithm produces a $(O(\phi \log n), \phi)$-ED by recursively partitioning $G$ along $\phi$-sparse cuts until no more such cuts are found.


\begin{algorithm}[H]
	\caption{\ed: recursive procedure for computing a $(O(\phi \log n), \phi)$-ED of $G$}
	\label{alg:metaed}
	
	\begin{algorithmic}[1]
		\LeftComment{$\phi \in (0,1)$ is the sparsity parameter}
		\Procedure{$\ed(G)$}{} \Comment{$G=(V,E)$ is the input graph}
		
		\If{$G$ is a $\phi$-expander} \Return $\{V\}$
		\Else
		\State $S \gets $ a $\phi$-sparse cut of $G$ \label{eq:sparsecut}
		\State \Return $\ed(G[S]) \cup \ed(G[V \setminus S])$
		\EndIf
		
		\EndProcedure
	\end{algorithmic}
\end{algorithm}
\noindent
Many algorithmic constructions of EDs are essentially efficient implementations of \cref{alg:metaed}. Adapting \cref{alg:metaed} to dynamic streams comes with its own set of challenges.
When the input is given in a dynamic stream, one can only afford to store a limited amount of information about the input graph. Since \cref{alg:metaed} only needs to measure the sparsity of cuts, it seems enough to have access to cut sizes and volumes. Both these quantities are preserved by cut sparsifiers. According to the classical definition~\cite{benczurkarger}, a $\delta$-cut sparsifier is a weighted subgraph $H=(V,E',w)$ of the input $G=(V,E)$, where for every cut $S \subseteq V$ one has
\begin{equation*}
	(1-\delta)\cdot \uncutg{S} \le \wcutg{S}\le (1+\delta) \cdot \uncutg{S} \, .
\end{equation*}
It is known that such a sparsifier can be constructed in dynamic streams using $\otil(n/\delta^2)$ bits of space~\cite{sparsdynamicstreams}. With the same space requirement we can also measure the sparsity of cuts up to a $(1\pm \delta)$ multiplicative error. Having this in mind, it is natural to consider the following algorithm: first, read the stream and construct a cut sparsifier, then run \cref{alg:metaed} on it and output the resulting clustering as the expander decomposition. However, this approach does not work as is. As it turns out, more information is needed about the graph than what is captured by the regular notion of a cut sparsifier.

The problem with this approach becomes immediately apparent when one considers the second recursion level. As the process recurses on the two sides of a sparse cut $S$, it will repeat the procedure on the subgraphs $G[S]$ and $G[V \setminus S]$. Unfortunately, the cut preservation property of the sparsifier does not carry over to those subgraphs, making it inadequate for estimating the sparsity of the cuts in those graphs. 
In fact, the notion of expander decomposition itself already operates on the subgraphs, as it is required that each cut in each cluster of the decomposition is expanding.

\subsubsection*{Testing expansion of subgraphs}
The reasoning from above gives rise to the following sketching problem: produce a sparsifier that can be used to check that any subgraph is a $\phi$-expander. To solve it, the authors of~\cite{streamexpdec} introduced the concept of a $(\delta,\psi)$-power-cut sparsifier. Its property is that, for any cluster $U \subseteq V$, with high probability all cuts $S \subseteq U$ are preserved within an additive-multiplicative error.
Recall that $\uncut{S}{U}$ is equal to the size of the cut $S$ inside the induced subgraph $G[U]$.
When talking about sparsifiers in this section, we will slightly abuse notation by assuming that there is some instance $H=(V,E',w)$ of it and denote by $\wcut{S}{U}$ the size of the cut $S$ in the subgraph $H[U]$ of this sparsifier~$H$. Then we can write the guarantee of a $(\delta, \psi)$-power-cut sparsifier as follows\footnote{A similar sparsifier construction was proposed by~\cite{sublinhierarchical}. Their construction has an additive error of $\psi |S|$, so the dependence is on the number of vertices instead of the volume}:
\begin{equation*}
	\forall \, S \subseteq U, \quad  (1-\delta)\cdot \uncut{S}{U}-\psi \cdot \vol(S) \le \wcut{S}{U} \le (1+\delta)\cdot \uncut{S}{U}+\psi \cdot \vol(S) \, .
\end{equation*}
The authors also give a dynamic stream construction, which uses $\otil(n/\delta\psi)$ bits of space by sampling edges proportionally to the degrees of their endpoints.


To check the $\phi$-expansion of subgraphs up to a small constant multiplicative error, it is enough to set $\delta$ to be a small constant. However, the multiplicative error parameter $\psi$ must be $\lesssim \phi$. 
This is a significant downside of this construction, as it was shown in \cite{streamexpdec} that a $(\delta, \phi)$-power-cut sparsifier must have at least $\Omega(n /\phi)$ edges. In fact, their lower bound is more general, and holds for general subgraphs\footnote{The lower bound instance is a $1/\phi$-regular graph. There, each edge by itself forms a $\phi$-expander, while any pair of vertices without an edge between them is not. Being able to test the $\phi$-expansion of the majority of those small subgraphs would imply being able to recover the majority of edges in the graph.}.
In other words, {$\Omega(n/\phi)$} space is necessary in order to test {$\phi$}-expansion of general subgraphs. Consequently, new tools must be used to have any hope of getting an algorithm independent~of~$1 / \phi$.

\paragraph*{Our contribution: sparsification of boundary-linked subgraphs}

As it turns out, solving the expansion testing problem, as it was stated, is not necessary. Recently, in the breakthrough work of \cite{boundlink}, it was shown that one could demand additional properties from the expander decomposition, and it will still exist at the price of increasing the number of inter-cluster edges by a small multiplicative factor.

More formally, for a set of vertices $U \subseteq V$ and $\tau > 0$, we denote by $G[U]^\tau$ an induced subgraph of $G$ on $U$, where additionally for each edge $e=\{u,v\}$ that connects a vertex $u$ inside $U$ with a vertex $v$ outside of $U$, $\tau$ self-loops are added to $u$. The graph $G[U]^\tau$ is called a $\tau$-boundary-linked subgraph. Note that if $G[U]^\tau$ is a $\phi$-expander, then so necessarily is $G[U]$, but not the other way around.


Then it is possible, for a given $\phi$, and $\eps = \otil(\phi), \, \tau \approx 1/\phi$, to construct an $(\eps, \phi)$-expander decomposition $\mathcal{U}$ where for each cluster $U \in \mathcal{U}$, $G[U]^\tau$ is a $ \widetilde{\Omega}(\phi)$-expander~\cite{boundlink}. Such a decomposition is called a boundary-linked expander decomposition.

A key observation is that if in our algorithm we were aiming to construct a boundary-linked expander decomposition, the testing problem would only involve checking that any given boundary-linked subgraph is $\phi$-expanding. Indeed, this problem is much easier than the original one and can be solved in space \smash{$\otil(n)$}. We will show how to do it in two steps: first, we will discuss how to strengthen the power-cut sparsifier, and then prove that this strengthening is enough to resolve the problem.

\paragraph*{Achieving additive error in the global cut}
As was noted in \cite{streamexpdec}, the idea behind constructing a power-cut sparsifier was to reanalyse the guarantee given by the construction of~\cite{stspectral} for sparsifying expanders. In other words, a power-cut sparsifier results from a more rigorous analysis of an existing sparsifier. The problem with this approach is that the sparsifier of~\cite{stspectral} is relatively weak to begin with: it only preserves cuts in expanders, while other constructions can preserve them in all graphs~\cite{benczurkarger,fungspars}.
To strengthen the guarantee, we give the same treatment to the construction of \cite{fungspars} in \cref{sec:spars}. For the sake of simplicity and to gain an intuition for why this kind of sparsification is at all possible, we discuss here how to do that with the classical construction of \cite{karger1994random} by closely following the original proof.

We show that, given a graph $G=(V, E)$ with a minimum cut of size $k$, it is possible to construct a sparsifier $H$ of $G$ such that every cut inside any given subgraph $G[U]$ of $G$ is preserved with high probability with the following guarantee:
\begin{equation}
	\label{eq:to-cut-guar}
	\forall \, S \subseteq U, \quad \uncut{S}{U} - \delta \cdot \uncutg{S} \le \wcut{S}{U}  \le \uncut{S}{U} + \delta \cdot \uncutg{S}\,.
\end{equation}
In this paper, a sparsifier with property~\eqref{eq:to-cut-guar} is called a cluster sparsifier (see \cref{def:spars}).
To achieve~\eqref{eq:to-cut-guar}, consider using the same process as in \cite{karger1994random}: sample each edge with the same probability $p \approx {\delta^{-2}/ k}$.

To see why this works, fix a subgraph $U \subseteq V$, and consider any cut $S$ in $U$. 
We wish to show that the size of the cut $S$ inside $U$ concentrates well after sampling. In the original proof, this is done by simply applying a Chernoff bound. In our case, this bound would look like this:
\begin{equation*}
	\Pr[|\wcut{S}{U} - \uncut{S}{U}| \geq \delta \cdot \uncut{S}{U}] \leq \exp\left(-\frac{1}{3} \delta^2 p \cdot \uncut{S}{U} \right) \, .
\end{equation*}
\noindent
However, this is insufficient, as the probability depends on the cut size inside $G[U]$. As we have no lower bound on its size, unlike with the sizes of global cuts, the second part of the argument of \cite{karger1994random} cannot be applied. Instead, we apply an additive-multiplicative version of the Chernoff bound (see for example \cite{streamexpdec}), that allows us to compare the approximation error with a bigger value than its expectation. This gives us
\begin{equation*}
	\Pr[|\wcut{S}{U} - \uncut{S}{U}| \geq \delta \cdot \uncutg{S}] \leq 2\exp\left(-\frac{1}{100} \delta^2 p \cdot \uncutg{S} \right) \, .
\end{equation*}
Expressing the probability in terms of the global cut allows us to use the cut counting lemma~\cite{karger1993global}, which bounds the number of global cuts of size at most $\alpha k$ by $n^{2\alpha}$,  for $\alpha \geq 1$. The proof is concluded by associating each global cut with a local cut in $G[U]$ and taking the union bound over them.

In order to get the guarantee of~\eqref{eq:to-cut-guar} for all graphs, not only those with a minimum cut of size $k$, one can sample edges proportionally to the inverse of their connectivity (as in the work of \cite{fungspars}) as opposed to uniformly. Moreover, one can implement such sampling scheme in dynamic streams, using the approach of Ahn, Guha, and McGregor~\cite{sparsdynamicstreams}. We defer the details of these reanalyses to \cref{sec:spars}, where we show how to construct cluster sparsifiers with property~\eqref{eq:to-cut-guar} in dynamic streams.

\paragraph*{Benefits of boundary-linked graphs}
To see why the cluster sparsifier is enough to solve the boundary-linked $\phi$-expansion testing problem, consider the following reasoning. Set the self-loop parameter $\tau$ equal to $b / \phi$, for some $b \gg \phi$ and $b \ll 1$. Fix a cluster $U$, for which $G[U]^{b / \phi}$ is an $\widetilde{\Omega}(\phi)$-expander. 
The crucial fact is that the size of any cut inside $G[U]^{b / \phi}$ is lower bounded by its size in the global graph up to a small polylogarithmic factor in the following way:
\begin{equation}
	\label{eq:to_cut_lb}
	\uncut{S}{U} \geq \widetilde{\Omega}(b) \cdot \uncutg{S} \, .
\end{equation}
We will now explain the derivation of the above equation in detail.
For a cut $\emptyset \neq S \subsetneq U$, we denote by $\unbor{S}{U} = \uncutg{S} - \uncut{S}{U}$ the number of edges going from $S$ to $V \setminus U$. Note that by the definition of~$G[U]^{b / \phi}$, the volume of any cut $S$ inside of it is equal to \smash{$\vol_{G[U]^{b / \phi}}(S) = \vol(S) + (\frac{b}{\phi} - 1)\unbor{S}{U}$}.
First, because we assume that $G[U]^{b / \phi}$ is an $\widetilde{\Omega}(\phi)$-expander, we have
\begin{equation*}
	\uncut{S}{U} \geq \widetilde{\Omega}(\phi)\cdot \vol_{G[U]^{b / \phi}}(S)\,.
\end{equation*}
Then, applying the aforementioned formula for $\vol_{G[U]^{b / \phi}}(S)$, we have
\begin{equation*}
	\uncut{S}{U} \geq \widetilde{\Omega}(\phi)\vol(S) + \widetilde{\Omega}(\phi)\left(\frac{b}{\phi} -1 \right)\unbor{S}{U} \,.
\end{equation*}
Since $b \gg \phi$ and dropping the first summand, the above simplifies to 
\begin{equation*}
	\uncut{S}{U} \geq  \widetilde{\Omega}(b)\unbor{S}{U} \,.
\end{equation*}
Finally, applying $\uncut{S}{U} + \unbor{S}{U} = \uncutg{S}$ and $b \ll 1$, we arrive back at equation~\eqref{eq:to_cut_lb}.

A consequence of equation~\eqref{eq:to_cut_lb} is that setting $\delta \approx 1/ b$ in the cluster sparsifier produces the following guarantee for $ \widetilde{\Omega}(\phi)$-expander subgraphs $G[U]^{b / \phi}$: 
\begin{equation*}
	\forall \, S \subseteq U, \quad \uncut{S}{U} - O(1) \cdot \uncut{S}{U} \le \wcut{S}{U}  \le \uncut{S}{U} + O(1) \cdot \uncut{S}{U}\,,
\end{equation*}
where the $O(1)$ can be made arbitrarily small. In other words, we achieve a multiplicative approximation guarantee on subgraphs of interest, which is enough to solve the testing problem. Since $b$ can be set to be $1/\pylog n$, the final sparsifier size does not depend on $1/\phi$. In this sense, such sparsifier can be thought of as a ``universal'' sketch of the graph that allows to solve the testing problem for all $\phi$ simultaneously.

Even though we now know how to solve the testing problem, an implementation of \cref{alg:metaed} would need to actually find a sparse cut when the test fails (and in particular, it needs to find a \textit{balanced} sparse cut, see the next section). In \cref{sec:balsparse}, we show that cluster sparsifiers enable to solve this harder task too: consider an offline algorithm that either correctly determines a graph to be a boundary-linked expander or finds a (balanced) sparse cut; we show that one can run such algorithm on a subgraph of the sparsifier as a black box and obtain essentially the same result as if the algorithm was run on the corresponding subgraph of the original graph. In this sense, cluster sparsifiers serve as small error proxies to the original graph for expander-vs-sparse-cut type of queries on vertex-induced subgraphs.

\subsubsection*{Low depth recursion}
Another problem that arises when using the recursive approach is that the subgraph sparsification guarantee of both power-cut sparsifiers and cluster sparsifiers is only probabilistic, and it can be shown that it cannot be made deterministic \cite{streamexpdec}. This means that after finding a sparse cut~$S$ in a subgraph $G[U]^{b / \phi}$ of a sparsifier, we cannot claim that the same sparsifier would preserve the cuts inside the new graph $G[S]^{b / \phi}$ with high probability, as it is dependent on another cut that we have already found inside the sparsifier. This means we cannot use the same sparsifier for two different calls to \cref{alg:metaed} inside the same execution path.
However, we can share a sparsifier among all the calls at the same recursion level since these will operate on independent portions of~$G$. Therefore, we want to have a separate sparsifier for every recursion level. This means that in order to minimize space requirements, it is crucial to have a small recursion depth.

To have small recursion depth, the algorithmic approach of \cite{streamexpdec,triang,ChangPSZ21} enforces the sparse cut $S$ from line~\eqref{eq:sparsecut} to be balanced, i.e., none of $S$ and $V\setminus S$ is much larger than the other. In particular, they only recurse on the two sides of a cut if $\vol(S) \gtrsim \epsilon \vol(V \setminus S)$. 
When there is no balanced sparse cut and yet the input is not an expander, a lemma of Spielman and Teng~\cite{stspectral} suggests there should be an $\Omega(\phi)$-expander $G[S']$ that accounts for a $(1-O(\epsilon))$-fraction of the total volume. An algorithmic version of this structural result allows us to iteratively trim off a small piece of the graph until such $S'$ is found. At this point, the algorithm of~\cite{streamexpdec} can simply return $\{S'\} \cup (\cup_{u \in V \setminus S'} \{u\})$ as an ED, with at most $O(\epsilon|E|)$ inter-cluster edges between singletons. As the volume of the cluster multiplicatively decreases by $1 - O(\eps)$ after each call, this gives recursion depth at most $\otil(1/\epsilon) \approx \otil(1/\phi)$.

\paragraph{Our contribution: adaptation of trimming.} We show that a simple refinement of this approach allows us to adapt the framework of \cite{offlineexpdec}, which leads to an algorithm with recursion depth independent of  $\phi$. We run the same algorithm, but instead of separating each vertex in $V \setminus S'$ into its own singleton cluster, we recurse with \cref{alg:metaed} on the whole set $V \setminus S'$. Conceptually, this implements an analogue of the trimming step of \cite{offlineexpdec}.

This means that at the end, fewer edges in $V \setminus S$ become inter-cluster edges. Because of that,
we can strengthen the balancedness requirement: we recurse on the two sides of a sparse cut only if \smash{$\vol(S) \gtrsim \frac{1}{C} \vol(V \setminus S)$} for a large constant $C$. This allows us to trim more vertices each time, resulting in $\otil(1)$ depth of the trimming step.  On the other hand, because in each call to \cref{alg:metaed} the volume of clusters passed to recursive calls is decreased by at least a constant factor, the total recursion depth becomes at most $\otil(1)$.

Other than the refinement discussed above, our space efficient implementation of \cref{alg:metaed}, as well as  the iterative procedure to find the large expander $S'$, are almost the same as the ones of~\cite{streamexpdec} (which in turn are  inspired by the one of~\cite{triang}). The details of the algorithms are given in \cref{sec:algexpdec}.

\subsubsection*{Putting it all together}
Combining the ideas illustrated in the two sections above, neither the sparsifier's size nor the recursion depth depend on $\phi$, thus giving a sparsity-independent space algorithm for boundary-linked expander decomposition  (BLD for short). This result is stated in terms of parameters $b,\epsilon,\phi \in (0,1), \, \gamma \ge 1$: a $(b, \epsilon, \phi, \gamma)$-BLD is a partition $\mathcal{U}$ with  at most an $\epsilon$ fraction of crossing edges and every $U \in \mathcal{U}$ induces a $\phi/\gamma$-expander $G[U]^{b/\phi}$ (see \cref{sec:prelims} and \cref{def:expdec}). Using this terminology, we obtain the following result.

\begin{restatable}[Exponential time decoding BLD]{theorem}{mainbld}
	\label{th:main}
	Let $G=(V, E)$ be a graph given in a dynamic stream, and let $b \in (0,1)$ be a parameter such that $b \le 1/\log^2 n$. Then, there is an algorithm that maintains a linear sketch of $G$ in $\otil(n/b^3)$ space. For any $\epsilon \in [n^{-2}, b \log n]$, the algorithm decodes the sketch to compute, with high probability and in $\otil(n/b^3)$ space and $2^{O(n)}$ time, a $(b, \epsilon, \phi, \gamma)$-BLD of $G$ for
	\begin{equation*}
		\phi = \Omega\left(\frac{\epsilon}{\log n}\right) \quad \text{ and } \quad \gamma = O(1) \, .
	\end{equation*}
\end{restatable}
\noindent
From this, one can easily conclude our main result, restated here for convenience of the reader.

\begin{mainthactual}[ED algorithm -- exponential time decoding]
	\label{cor:maingeneral}
	Let $G=(V, E)$ be a graph given in a dynamic stream. Then, there is an algorithm that maintains a linear sketch of $G$ in $\otil(n)$ space. For any  \smash{$\phi \in (0, 1)$} such that $ \phi \le c/\log^2 n$ for a small enough constant $c>0$, the algorithm decodes the sketch to compute a $(O(\phi \log n), \phi)$-ED of $G$ with high probability, in $\otil(n)$ space and $2^{O(n)}$ time.
\end{mainthactual}
\noindent
We remark that for \smash{$\phi \ge \Omega(1/\log^2n)$}, one can use the algorithm of~\cite{streamexpdec} to still have an $\otil(n)$ space construction of a $(O(\phi \log n), \phi)$-ED.
\begin{proof}
	First note that without loss of generality we can assume $\phi \ge 1/n^2$, otherwise an ED can simply consists of the connected components of $G$ (which can be computed in dynamic streams in~$\otil(n)$ space~\cite{linmeas}). Then, we note that since $c$ is small enough and $\phi \le c / \log^2 n$, one can always define $\epsilon = C \cdot \phi \cdot \log n$ for an appropriate constant $C>0$ while ensuring $1/n^2 \le \epsilon \le 1/\log n$. We can thus prove the theorem by equivalently showing that there is an algorithm that maintains a linear sketch of $G$ in $\otil(n)$ space, and that for all \smash{$\epsilon \in [{1}/{n^2},{1}/{\log n}]$} decodes the sketch to compute, with high probability, an $(\epsilon, \Omega(\epsilon / \log n)))$-ED of $G$ in $\otil(n)$ space and $2^{O(n)}$ time.
	
	Let then \smash{$\epsilon \in [{1}/{n^2},{1}/{\log n}]$}. We use the algorithm from \cref{th:main} with parameter $\epsilon$ and a parameter $b$ of our choice. We need to meet two preconditions: $b \le 1/\log^2 n$ and $\epsilon \le b \log n$. Since we assume $\epsilon \le 1/\log n$, we can set $b=1/\log^2 n$, and all the prerequisites are fulfilled. Then the algorithm from \cref{th:main} runs in $2^{O(n)}$ time and takes $\otil(n/b^3)=\otil(n)$ bits of space. The output $\mathcal{U}$ is a $(b, \epsilon, \phi, \gamma)$-BLD of $G$ with high probability, where $\phi = \Omega(\epsilon/\log n)$ and $\gamma =O(1)$. Since a $(b, \epsilon, \phi, \gamma)$-BLD of $G$ is an $(\epsilon,\phi/\gamma)$-ED of $G$, we have obtained an $(\epsilon, \Omega(\epsilon/\log n))$-ED of $G$ with high probability.
\end{proof}
\noindent
The exponential time in the decoding is due to the subtask of finding a balanced sparse cut. As we show, one can make the decoding time polynomial by resorting to known offline approximation algorithms~\cite{offlineexpdec,liexpdec}. However, we only have $\log^{\Omega(1)} n$-approximations for finding a balanced sparse cut, and in particular, we do not expect (under NP-hardness and the Unique Games Conjecture) there to be a polynomial time $O(1)$-approximation~\cite{hardsparsecut}. Such super-constant factor error incurs some loss in the quality of decomposition and space requirement, which, nevertheless, remains independent of the sparsity.


\begin{restatable}[Polynomial time decoding BLD]{theorem}{mainbldpoly}
	\label{th:mainpoly}
	Let $G=(V, E)$ be a graph given in a dynamic stream, and let $b \in (0,1)$ be a parameter such that $b \le 1/\log^5 n$. Then, there is an algorithm that maintains a linear sketch of $G$ in \smash{$n/b^3 \cdot \log^{O(\log n / \log\frac{1}{b})} n $} space. For any $\epsilon \in [n^{-2}, b \log n]$, the algorithm decodes the sketch to compute, with high probability and in \smash{$n/b^3 \cdot \log^{O(\log n / \log\frac{1}{b})} n $} space and $\poly(n)$ time, a $(b, \epsilon, \phi, \gamma)$-BLD of $G$ for
	\begin{equation*}
		\phi = \Omega\left(\frac{\epsilon}{\log^4 n}\right) \quad \text{ and } \quad \gamma = \log^{O\left(\frac{\log n}{\log 1/b}\right)} n\, .
	\end{equation*}
\end{restatable}
\noindent
\cref{cor:mainpolygeneral}, restated here for convenience, then follows from \cref{th:mainpoly}.

\begin{mainthactual}[ED algorithm -- polynomial time decoding]
	\label{cor:mainpolygeneral}
	Let $G=(V, E)$ be a graph given in a dynamic stream, and let $b \in (0,1)$ be a parameter such that $b \le 1/\log^5 n$. Then, there is an algorithm that maintains a linear sketch of $G$ in \smash{$n/b^3 \cdot \log^{O(\log n / \log\frac{1}{b})} n $} space. For any {$\phi \in (0,1)$} such that \smash{$\phi \le b/\log^{C \cdot \log n / \log \frac{1}{b}}n$} for a large enough constant $C>0$, the algorithm decodes the sketch to compute a \smash{$(\phi \cdot \log^{O(\log n / \log\frac{1}{b})} n, \phi)$}-ED of $G$  with high probability, in \smash{$n/b^3 \cdot \log^{O(\log n / \log\frac{1}{b})} n $} space and~\smash{$\poly(n)$} time.
\end{mainthactual}
\noindent
We remark that setting, say, $b = 2^{-\sqrt{\log n}}$  in the above result gives a $n^{1+o(1)}$ space algorithm for computing a $(\phi \cdot n^{o(1)}, \phi)$-ED for any \smash{$\phi \le 2^{-2C\log\log n \sqrt{\log n}}$}. For larger values of $\phi$, one can use the polynomial time algorithm of ~\cite{streamexpdec} to still get a $n^{1+o(1)}$ space construction for a  $(\phi \cdot n^{o(1)}, \phi)$-ED.
\begin{proof}
	As in the proof of \cref{cor:maingeneral}, we can assume $\phi \ge 1/n^2$. Also, by virtue of $C$ being a large enough constant and  \smash{$\phi \le b/\log^{C \cdot \log n / \log \frac{1}{b}}n$}, one can always define \smash{$\epsilon = \phi \cdot \log^{C \cdot \log n / \log \frac{1}{b}}n$} while ensuring $n^{-2} \le \epsilon \le b \log n$. Then, we equivalently prove that for any $b \in (0,1)$ with $b \le 1/\log^5 n$ there is an algorithm that maintains a linear sketch of~$G$ in \smash{$n/b^3 \cdot \log^{O(\log n / \log\frac{1}{b})} n $} space, and that for any $\epsilon \in (0,1)$ such that $n^{-2} \le \epsilon \le b \log n$ decodes the sketch to compute, with high probability, an \smash{$(\epsilon, \epsilon/\log^{O(\log n / \log\frac{1}{b})} n)$}-ED of~$G$ in \smash{$n/b^3 \cdot \log^{O(\log n / \log\frac{1}{b})} n $} space and \smash{$\poly(n)$} time.
	
	Let then $b,\epsilon \in (0,1)$ with $b \le 1/\log^5 n$ and $n^{-2} \le \epsilon \le b \log n$. We use the algorithm from \cref{th:mainpoly} with the same parameters $b$ and $\epsilon$, since every admissible pair of parameters $b$ and $\epsilon$ fulfils the conditions of \cref{th:mainpoly}. The space complexity is also the same, and the running time is $\poly(n)$. Again, observe that a $(b, \epsilon, \phi, \gamma)$-BLD of $G$ is an $(\epsilon,\phi/\gamma)$-ED of $G$. Hence the claim, since \cref{th:mainpoly} gives
	\begin{equation*}
		\frac{\phi}{\gamma} = \Omega\left(\frac{\epsilon}{\log^4 n}\right) \cdot \frac{1}{\log^{O\left(\frac{\log n}{\log 1/b}\right)} n} = \frac{\epsilon}{\log^{O\left(\frac{\log n}{\log 1/b}\right)} n} \, .
	\end{equation*}
\end{proof}

\noindent
The proofs of \cref{th:main} and \cref{th:mainpoly} can be found in \cref{sec:algexpdec}.
